\documentclass[11pt]{scrartcl}
\usepackage[utf8]{inputenc}
\usepackage[sexy]{evanchen} 
\usepackage{graphicx}
\graphicspath{{8 Outsource K Math/Gambar/}}

\title{Quadratic Residue}
\author{Azzam and Akeyla}

\begin{document}

\maketitle

\section{Materi pengantar}
\subsection{Order}
Misalkan $a$ dan $m$ adalah bilangan bulat dengan $m \neq 0$ dan $m \neq 1$ serta $FPB(a,m) = 1$. Suatu bilangan asli $x$ disebut order dari $a$ dalam modulo $m$ jika $$a^x \equiv 1 \pmod{m}$$ dan untuk setiap bilangan asli yang kurang dari $x$, $a^x \not\equiv 1 \pmod{m}$. Hal ini biasa dinotasikan dengan $x = ord_m(a)$
\subsubsection{Sebuah teorema yang berguna}
\begin{theorem}
    Jika $z$ adalah bilangan asli sedemikian hingga $a^z \equiv 1 (mod m)$, maka $ord_m(a) | z$.
\end{theorem}
\textbf{Bukti :} Gunakan algoritma pembagian.\\
Sebagai akibat dari teorema ini, $ord_m(a) | \phi(m)$


\subsection{Akar Primitif}
Diberikan sebuah bilangan asli $m$. Sebuah bilangan bulat $g \not\equiv \pmod{m}$ disebut sebagai akar primitif modulo $p$ apabila $ord_m(g) = \varphi(m)$.
\subsubsection{Properti yang berguna}
Jika $g$ adalah akar primitif dalam modulo $p$, maka \[\{g^1, g^2, ..., g^{\varphi(m)}\} \pmod{m}\] memuat semua kemungkinan residu modulo $m$ yang relatif prima dengan $m$

\section{Quadratic Residue}
Sebuah bilangan bulat $a$ disebut sebagai quadratic residue (QR) dalam modulo $n$ jika dan hanya jika ada sebuah bilangan bulat $x$ yang memenuhi $x^2 \equiv a \pmod{n}$. Jika tidak ada nilai $x$ yang memenuhi, $a$ disebut sebagai quadratic non-residue (QNR).
\subsection{Notasi Legendre}
Untuk setiap bilangan prima $p$ dan bilangan bulat $x$, definisikan notasi Legendre sebagai:
\[\left(\dfrac{x}{p}\right) = 
\begin{cases}
    1,  \text{jika $x$ adalah QR},\\
    0, \smallskip \text{jika $x \equiv 0 \pmod{p}$},\\
    -1, \smallskip \text{selain itu}
\end{cases}\]
\subsection{Beberapa properti QR}
\begin{enumerate}
    \item Untuk setiap bilangan prima ganjil $p$, tepat setengah dari residu modulo $p$ merupakan QR dan setengahnya lagi merupakan QNR
    \item Jika $a, b$ adalah bilangan bulat dan $p$ adalah bilangan prima, maka \[\left(\dfrac{a}{p}\right)\left(\dfrac{b}{p}\right) = \left(\dfrac{ab}{p}\right)\]
    \item Jika $a \equiv b \pmod{p}$ maka \[\left(\dfrac{a}{p}\right) = \left(\dfrac{b}{p}\right)\]
    \item (Kriteria Euler) Jika $a$ adalah bilangan bulat dan $a$ adalah bilangan prima maka \[\left(\dfrac{a}{p}\right) \equiv a^{\frac{p - 1}{2}} \pmod{p}\]
    \item Jika $p$ adalah bilangan prima ganjil maka \[\left(\dfrac{2}{p}\right) = (-1)^{\frac{p^2 - 1}{8}}\]
    \item (Gauss's Law of Quadratic Reciprocity) Jika $p, q$ merupakan dua bilangan prima ganjil berbeda, maka \[\left(\dfrac{p}{q}\right)\left(\dfrac{q}{p}\right) =  (-1)^{\frac{p - 1}{2} \cdot \frac{q - 1}{2}}\]
\end{enumerate}


\section{Contoh Soal}
\begin{example}
    Diberikan bilangan bulat $x$ dan $y$ yang memenuhi $5|x^2 + xy + y^2$. Buktikan bahwa $25$ juga habis membagi $x^2 + xy + y^2$.
\end{example}
\begin{proof}
    Jika $5|x$ maka $5$ juga harus habis membagi $y$ dan sebaliknya. Akibatnya, $5^2 = 25$ habis membagi $x^2 + xy + y^2$. Jadi, asumsikan bahwa $5$ tidak habis membagi kedua bilangan ini. Amati bahwa $5|4(x^2 + xy + y^2) = (2x + y)^2 + 3y^2$. Akibatnya, $(2x + y)^2 \equiv -3y^2 \pmod{5}$. Dengan mengalikan kedua ruas dengan invers $y$ dalam mod $5$, didapat $-3 \equiv (2xy^{-1} + 1)^2 \pmod{5}$. Jadi, $-3$ adalah QR dalam mod $5$. Oleh karena itu, 
    \begin{align*}
        1 &= \left(\dfrac{-3}{5}\right) \\
        &\equiv (-3)^{\frac{5 - 1}{2}} \pmod{5}\\
        &\equiv 9 \pmod{5}
    \end{align*}
    yang jelas salah. Jadi, asusmsi salah dan haruslah $5$ membagi kedua bilangan tersebut. Akibatnya, $25$ membagi $x^2 + xy + y^2$
\end{proof}

\begin{example}
    Apakah $37$ adalah QR dalam modulo $83$?
\end{example}

\begin{proof}
    Dengan menggunakan quadratic reciprocity, didapat \[\left(\dfrac{37}{83}\right)\left(\dfrac{83}{37}\right) = (-1)^{18\cdot 41} = 1\] Lebih lanjut, karena $\left(\dfrac{83}{37}\right) = \left(\dfrac{9}{37}\right) = 1$, maka $\left(\dfrac{37}{83}\right) = 1$ dan $37$ adalah QR dalam modulo $83$.
\end{proof}

\section{Latihan}
\begin{problem}
    Kapankah $-1$ merupakan QR dari sebuah bilangan prima $p$?
\end{problem}

\begin{problem}
    Jika $p$ adalah bilangan prima ganjil dan terdapat bilangan bulat $x, y$ yang tidak habis dibagi $p$ yang memenuhi $p|x^2 + y^2$ maka $p \equiv 1 \pmod{4}$
\end{problem}

\begin{problem}
    Buktikan bahwa tidak ada solusi bilangan asli dari persamaan $4ab - a - b = c^2$
\end{problem}

\begin{problem}
    (Vietnam TST 2004) Buktikan bahwa $2^n + 1$ tidak punya faktor prima yang berbentuk $8k - 1$
\end{problem}

\begin{problem}
    Diberikan sebuah bilangan prima $p > 3$. Definisikan $S = \{0^3, 1^3, 2^3, ..., (p - 1)^3\}$. Buktikan bahwa $S$ merupakan kelas residu lengkap modulo $p$ jika dan hanya jika $p \equiv 2 \pmod{3}$.
\end{problem}

\begin{problem}
    (Taiwan MO 1997) Misal $k = 2^{2^n} + 1$ untuk sebuah bilangan asli $n$. Buktikan bahwa $k$ merupakan bilangan prima jika dan hanya jika $k$ adalah faktor dari $3^{\frac{k - 1}{2}} + 1$
\end{problem}

\begin{problem}
    Tentukan faktor prima terkecil dari $12^{2^{15}} + 1$
\end{problem}

\begin{problem}
    (American Mathematical Monthly) Cari semua bilangan asli $n$ yang memenuhi $2^n - 1|3^n - 1$
\end{problem}

\begin{problem}
    (Polandia MO 2019) Misal $p$ adalah sebuah bilangan prima dan $r$ adalah bilangan asli yang memenuhi $p|r^7 - 1$. Buktikan bahwa apabila terdapat bilangan bulat $a, b$ yang memenuhi $p|r + 1 - a^2$ dan $p|r^2 + 1 - b^2$, maka terdapat sebuah bilangan bulat $c$ yang memenuhi $p|r^3 + 1 - c^2$.
\end{problem}

\newpage
\section{Solusi}
\begin{problem}
    Tentunya, jika $p = 2$, $-1$ adalah QR. Untuk $p$ ganjil, gunakan Kriteria Euler, dan didapat bahwa $p \equiv 1 \pmod{4}$ agar $-1$ adalah QR dalam mod $p$
\end{problem}

\begin{problem}
    Perhatikan bahwa $x^2 \equiv -y^2 \pmod{p}$ dan \[\left(\dfrac{x^2}{p}\right) = \left(\dfrac{-y^2}{p}\right) = \left(\dfrac{-1}{p}\right)\left(\dfrac{y^2}{p}\right) = \left(\dfrac{-1}{p}\right)\]
    Persamaan terakhir dikarenakan $y^2$ pasti merupakan QR dalam mod $p$. Selain itu, karena $x^2$ juga QR dalam mod $p$, maka $-1$ juga QR dalam mod $p$. Dengan soal pertama, dapat disimpulkan bahwa $p \equiv 1 \pmod{4}$
\end{problem}

\begin{problem}
    Kalikan persamaan tersebut dengan $4$ kemudian tambahkan dengan $1$, didapat $(4a - 1)(4b - 1) = 4c^2 + 1$. Ruas kiri memiliki setidaknya satu faktor prima berbentuk $-1 \pmod{4}$. Padahal, dari soal nomor $2$, semua faktor dari ruas kanan harus berbentuk $1 \pmod{4}$.
\end{problem}

\begin{problem}
    Ambil faktor prima $p$ dari bilangan $2^n + 1$ tersebut. Jika $n$ genap maka $-1$ adalah QR dalam mod $p$ dan akibatnya $p \equiv 1 \pmod{4}$ dan tidak mungkin $-1 \pmod{8}$. Jika $n$ ganjil, maka $2^{n + 1} \equiv -2 \pmod{p}$ yang membuat $-2$ adalah QR mod $p$. Namun, \[\left(\dfrac{-2}{p}\right) = \left(\dfrac{-1}{p}\right)\left(\dfrac{2}{p}\right) = (-1)^{\frac{p - 1}{2} + \frac{p^2 - 1}{8}}\] yang membuat $p \equiv -1 \pmod{8}$ tidak memungkinkan.
\end{problem}

\begin{problem}
    Asumsikan ada dua residu berbeda $m, n$ yang memenuhi $m^3 \equiv n^3 \pmod{p}$. Dikarenakan $m - n \not\equiv \pmod{p}$, didapat $m^2 + mn + n^2 \equiv 0 \pmod{p}$. Manipulasi bentuk tersebut (seperti pada contoh 1) akan menghasilkan $(2m + n)^2 \equiv -3n^2 \pmod{p}$ yang membuat $-3$ adalah QR dalam mod $p$. Dapat dicek dengan quadratic reciprocity bahwa hal ini hanya dapat terjadi saat $p \equiv 2\pmod{3}$.
\end{problem}

\begin{problem}
    Misal $k$ adalah faktor dari $3^{(k - 1)/2} + 1$. Maka, $z = ord_k(3)$ tidak akan membagi $(k - 1)/2 = 2^{2^{n - 1}}$ tetapi akan membagi $k - 1 = 2^{2^n}$ Artinya, $z = 2^{2^n} = k - 1$. Tentunya, $ord_k(3) = k - 1$ menunjukkan bahwa $k$ adalah bilangan prima (jika tidak, maka $\varphi(k) < k - 1$ padahal $ord_k(3)|\varphi(k)$)
    Sebaliknya, jika $k$ prima, maka \[3^{\frac{k-1}2}\equiv \left(\frac3k\right)=(-1)^{\frac{(3-1)(k-1)}4}\left(\frac k3\right)=\left(\frac k3\right)=\left(\frac 23\right)=-1\pmod k,\] karena $k = 2^{2^n} + 1\equiv 1 \pmod{4}$ dan $\equiv 2 \pmod{3}$.
\end{problem}

\begin{problem}
    Misal $p$ adalah faktor prima tersebut. Dikarenakan $p|12^{2^{16}} - 1$, maka $ord_p(12)|2^{16}$. Namun, karena $12^{2^{15}} \not\equiv 1 \pmod{p}$, didapat $ord_p(12)\not|2^{15}$ dan akibatnya, $ord_p(12) = 2^{16}$. Jadi, $2^{16}|\varphi(p) = p - 1$ dan $p \geq 2^{16} + 1 = 65537$. \textbf{Dapat dicek bahwa $65537$ prima}, jadi cukup cek apakah $q = 65537$ merupakan faktor dari $12^{2^{15}} + 1$ atau tidak. Padahal,
    \[12^{2^{15}} + 1 = 2^{q - 1}\cdot 3^{(q - 1)/2} + 1 \equiv 3^{(q - 1)/2} + 1 \pmod{q}\] Jadi, cukup dicek apakah $3$ merupakan QNR dalam mod $q$. Dengan quadratic reciprocity law, didapat \[\left(\dfrac{3}{q}\right) = (-1)^{\frac{(3 - 1)(q - 1)}{4}}\left(\dfrac{q}{3}\right) = \left(\dfrac{q}{3}\right) = \left(\dfrac{2}{3}\right) = -1\] dan dapat disimpulkan bahwa $q = 2^{16} + 1 = 65537$ memang faktor prima terkecil dari $12^{2^{15}}$.
\end{problem}

\begin{problem}
    $n = 1$ jelas memenuhi. Asumsikan $n > 1$. Jelas bahwa $n$ ganjil (kalau tidak, $2^n \equiv 1 \pmod{3})$. Sehingga, $2^n - 1 \equiv 7 \pmod{12}$. Dikarenakan faktor prima ganjil pasti berbentuk $12k \pm 1$ atau $12k \pm 5$, maka $2^n - 1$ pasti memiliki setidaknya satu faktor prima berbentuk $12k \pm 5$, misalkan saja $p$. Dikarenakan $3^n - 1 \equiv 0 \pmod{p}$ dan $n$ ganjil, maka haruslah $3$ merupakan QR dalam mod $p$. Dengan quadratic reciprocity law, didapat $\left(\dfrac{p}{3}\right) = (-1)^{(p - 1)/2}$. Padahal, $\left(\dfrac{p}{3}\right) = \left(\dfrac{\pm 2}{3}\right) = \mp 1$. Kedua nilai ini tidak akan bisa sama. Jadi, tidak ada nilai $p$ yang mungkin dan tidak ada nilai $n > 1$ yang mungkin.
\end{problem}

\begin{problem}
    
\end{problem}

\end{document}